% =====================================================================
%  第 5 节：技术路线
% =====================================================================
\section{技术路线}
\label{sec:roadmap}

本节给出 CRP 循环的整体设计（\S\ref{sec:r-arch}）、四个模块的具体职责（\S\ref{sec:r-modules}）、以及可审计轨迹的设计（\S\ref{sec:r-trace}）。本节描述\textbf{计划}，实验结果将在后续版本中补充。

\subsection{整体架构}
\label{sec:r-arch}

\begin{figure}[htbp]
\centering
\begin{tikzpicture}[
  x=1cm, y=1cm,
  box/.style={draw=accent!70, fill=accentlight, rounded corners=2pt,
              minimum width=3.2cm, minimum height=0.95cm, align=center,
              font=\small, inner sep=3pt},
  llmbox/.style={box, fill=orange!12, draw=orange!60!black},
  statbox/.style={box, fill=green!10, draw=green!45!black},
  arr/.style={-{Stealth[length=2.0mm]}, thick, accent!85},
  fb/.style={-{Stealth[length=2.0mm]}, thick, dashed, orange!75!black},
  lbl/.style={font=\scriptsize\itshape, gray}
]

% ---- 第 1 列：猜想 ----
\node[llmbox] (llm)  at (0,0)     {LLM 提议形式\\[-1pt]{\scriptsize 读样本摘要 + 已知定理}};
\node[box]    (lp)   at (0,-1.55) {线性规划定系数\\[-1pt]{\scriptsize 可行性 + 整数化}};
\node[box]    (cand) at (0,-3.10) {候选猜想集 $\mathcal{C}$\\[-1pt]{\scriptsize 形式 + 系数 + 紧性}};

% ---- 第 2 列：反驳 ----
\node[box] (search) at (4.9,0)     {反例搜索\\[-1pt]{\scriptsize 随机/结构/局部/精确}};
\node[box] (budget) at (4.9,-1.55) {预算分配\\[-1pt]{\scriptsize Beta–Bernoulli 汤普森}};

% ---- 第 3 列：筛选 ----
\node[statbox] (margin) at (9.8,0)     {违反裕度 $T_j$\\[-1pt]{\scriptsize 验证集上最大违反量}};
\node[statbox] (perm)   at (9.8,-1.55) {置换零分布\\[-1pt]{\scriptsize $\to$ $p$ 值}};
\node[statbox] (bh)     at (9.8,-3.10) {BH / BY 筛选\\[-1pt]{\scriptsize FDR $\le q$}};

% ---- 阶段标注 ----
\node[font=\small\bfseries, accent] at (0,0.95)     {(1)\ 猜想};
\node[font=\small\bfseries, accent] at (4.9,0.95)   {(2)\ 反驳};
\node[font=\small\bfseries, accent] at (9.8,0.95)   {(3)\ 筛选};

% ---- 数据流 ----
\draw[arr] (llm) -- (lp);
\draw[arr] (lp) -- (cand);
\draw[arr] (cand.east) -- ++(0.55,0) |- (search.west);
\draw[arr] (budget.north) -- (search.south);
\draw[arr] (search) -- (margin);
\draw[arr] (margin) -- (perm);
\draw[arr] (perm) -- (bh);

% ---- 反馈回路（绕行图外，避免穿越方框）----
\draw[fb] (bh.south) -- (9.8,-4.45) -- (-2.45,-4.45) -- (-2.45,0) -- (llm.west);
\node[lbl, anchor=north] at (4.0,-4.55) {失败原因 + 反例结构 $\Rightarrow$ 修正下一轮的形式提议};

\end{tikzpicture}
\caption{CRP 循环的架构。实线为数据流，虚线为反馈。左侧“猜想”与右侧“筛选”通过反馈回路耦合：筛选阶段暴露的失败模式被送回 LLM，用于修正下一轮的形式提议。}
\label{fig:arch}
\end{figure}

图 \ref{fig:arch} 的三列对应问题 \ref{prob:main} 的三个子问题。值得注意的是\textbf{反馈回路}：筛选不只是终点，它产生的“为什么这条候选不可信”的信息会回馈到生成端。这是系统被称为“智能体”而非“流水线”的实质原因——它有一个闭合的、基于结果调整策略的循环。

\subsection{模块职责}
\label{sec:r-modules}

\subsubsection{模块一：图样本与不变量计算}

\textbf{输入}：$n$ 的上界 $N_{\max}$、采样策略。
\textbf{输出}：不变量矩阵 $\bm{F}\in\R^{|\Sn|\times d}$ 及对应的图（以边表存储）。

关键决策是\textbf{如何处理 NP-难的不变量}。我们采取三级策略：

\begin{enumerate}[label=(\arabic*)]
  \item \textbf{$n\le 8$ 精确计算。}$|\G_8|=12{,}346$，用带剪枝的精确算法（$\alpha$ 用分支定界，$\chi$ 用 DSATUR + 回溯）可在分钟级完成。
  \item \textbf{$n=9$ 用界代替。}$|\G_9|=274{,}668$，精确计算色数不可行。改用区间 $[\omega,\ \Delta+1]$，并记录区间宽度作为“不确定度”。$\alpha$ 用精确算法（9 点图的独立数计算很快）。
  \item \textbf{$n\ge 10$ 采样。}用 $G(n,p)$（$p$ 扫描）、正则图、结构化图三类采样，每类若干千张。
\end{enumerate}

\begin{notebox}[一个需要报告的取舍]
用区间代替精确色数，会使“$\chi$ 参与的候选”的判定变得保守：若 $\chi$ 的真值在区间内变动能改变候选的真假，则该候选被标记为“不确定”而非“成立”或“被推翻”。这会降低召回率，但避免了假阳性。这个取舍将在实验部分量化。
\end{notebox}

\subsubsection{模块二：猜想生成}

\textbf{输入}：不变量矩阵、已知定理清单。
\textbf{输出}：候选集合 $\mathcal{C}$，每条含形式、系数、在发现集上的紧性。

流程：
\begin{enumerate}[label=\arabic*.]
  \item 计算样本摘要：各不变量的取值范围、分位数、极值图、两两相关系数；
  \item 把摘要与已知定理清单送入 LLM，要求其输出若干候选形式（JSON 格式，含不变量子集、方向、是否含取整）；
  \item 对每个线性形式，求解 LP \eqref{eq:lp}；不可行则记录并跳过；
  \item 对可行的解，做整数化（把系数缩放到小整数，检验是否仍可行）；
  \item 计算每条候选在发现集上的\textbf{紧性} $\max_G \bm{c}^\top\bm{f}(G)$，紧性接近 $0$ 的候选优先进入反驳。
\end{enumerate}

\subsubsection{模块三：反例搜索}

\textbf{输入}：候选 $\bm{c}$、预算 $b$。
\textbf{输出}：反例（若找到）或“未找到”及搜索充分性指标。

四种策略按预算比例混合执行，比例由上一轮的成败调整：

\begin{table}[H]
\centering
\caption{反例搜索的四种策略及其覆盖的反例形态。}
\label{tab:search}
\small
\begin{tabularx}{\linewidth}{@{}p{1.9cm} >{\raggedright\arraybackslash}X >{\raggedright\arraybackslash}p{4.3cm}@{}}
\toprule
策略 & 做法 & 擅长发现的反例 \\
\midrule
随机 & $G(n,p)$，$p\in\{0.1,\dots,0.9\}$，每个 $n$ 若干样本 & “典型”反例，$p$ 接近临界值时 \\
结构化 & 正则图、二部图、Mycielski 构造、Kneser 图、冠图、完全多部图 & 极端反例（$\chi$ 大而 $\omega$ 小等） \\
局部 & 从当前最优图出发做边翻转，爬山或模拟退火 & 精化：把违反量从接近 $0$ 推到正 \\
精确 & $n\le 8$ 时全枚举 & \textbf{确认}小图上无反例 \\
\bottomrule
\end{tabularx}
\end{table}

\subsubsection{模块四：统计筛选}

\textbf{输入}：存活候选及其在各集合上的违反裕度。
\textbf{输出}：保留集合、每条候选的 $p$ 值与 FDR 估计。

按 \S\ref{sec:f-screen} 的方案执行，另加两项：
\begin{itemize}
  \item \textbf{校准}：用已知定理（应被保留）与已知反例（应被推翻）构成的有标签子集，检验筛选的召回率与精确率；
  \item \textbf{敏感性分析}：对 $q$、置换次数 $B$、预算分配策略做网格扫描，报告结果对超参数的依赖。
\end{itemize}

\subsection{可审计轨迹}
\label{sec:r-trace}

整个循环的每一步写入 JSONL 文件，每行一条记录，字段包括：

\begin{table}[H]
\centering
\caption{轨迹记录的主要字段。}
\label{tab:trace}
\small
\begin{tabular}{@{}llp{7.2cm}@{}}
\toprule
字段 & 类型 & 含义 \\
\midrule
\texttt{round} & int & 循环轮次 \\
\texttt{stage} & str & \texttt{propose} / \texttt{refute} / \texttt{screen} \\
\texttt{candidate\_id} & str & 候选标识（形式的哈希） \\
\texttt{form} & object & 形式：不变量子集、方向、是否含取整 \\
\texttt{coeffs} & list & LP 求得的系数 \\
\texttt{llm\_input} & str & 送入 LLM 的完整提示 \\
\texttt{llm\_output} & str & LLM 的原始输出 \\
\texttt{search\_strategy} & str & 本轮使用的搜索策略 \\
\texttt{search\_budget} & int & 消耗的预算（评估次数） \\
\texttt{counterexample} & object & 反例的边表与不变量值（若有） \\
\texttt{margin} & float & 违反裕度 $T_j$ \\
\texttt{p\_value} & float & 置换检验得到的 $p$ 值 \\
\texttt{decision} & str & \texttt{keep} / \texttt{drop} / \texttt{uncertain} \\
\bottomrule
\end{tabular}
\end{table}

轨迹的作用有三：
\begin{enumerate}[label=(\arabic*)]
  \item \textbf{可复现}：给定轨迹，可重放整个循环（LLM 输出已记录，不需重新调用）；
  \item \textbf{可诊断}：分析“哪些形式的候选最容易被推翻”“哪类搜索策略最有效”；
  \item \textbf{可汇报}：中期与最终报告的图表直接由轨迹生成，避免手工整理。
\end{enumerate}

\subsection{验证计划}
\label{sec:r-valid}

由于“是否真的成立”无法完全判定，验证分三层：

\begin{enumerate}[label=(\roman*)]
  \item \textbf{正确性检验}：系统能否重现一批已知定理（Brooks、$\chi+\alpha\le n+1$、$\alpha\ge n/(\Delta+1)$ 等）？能否对已知的反例图（如 Grötzsch 图对某些断言）正确判定？
  \item \textbf{筛选有效性}：在有标签子集上，BH 筛选的实际 FDR 是否接近名义水平 $q$？与不做筛选（全保留）相比，精确率提升多少？
  \item \textbf{消融}：LLM 提议 vs.\ 随机提议 vs.\ 穷举提议，在相同预算下的产出质量差异；有预算分配 vs.\ 均匀分配。
\end{enumerate}

第 (i) 层是“系统能不能用”，第 (ii) 层是“统计方法有没有用”，第 (iii) 层是“智能体成分有没有用”。三层都通过，才能支持“本课题提出了一个有效方法”的主张；任一层失败，也是有价值的负面结果。

\subsection{当前进度与后续安排}
\label{sec:r-progress}

\begin{table}[H]
\centering
\caption{工作安排。表中“状态”截至本报告撰写时。}
\label{tab:schedule}
\small
\begin{tabular}{@{}clp{6.4cm}l@{}}
\toprule
阶段 & 内容 & 产出 & 状态 \\
\midrule
D0 & 选题论证、报告框架 & 本报告第一部分 & 完成 \\
D1 & 图枚举与不变量计算 & 不变量矩阵（$n\le 9$） & 待办 \\
D2 & LP 猜想生成器 & 首批候选 & 待办 \\
D3 & 反例搜索器 & 首轮反驳结果 & 待办 \\
D4 & 统计筛选与校准 & FDR 分析 & 待办 \\
D5 & 案例研究与消融 & 三条候选的深度分析 & 待办 \\
D6 & 中期汇报材料 & 幻灯片 & 待办 \\
\bottomrule
\end{tabular}
\end{table}
